% Same code as the original growth problem; qualitative solutions for Fall 2026.
\begin{blocksection}
\question Classify each function's work as linear, sublinear, or superlinear
relative to \lstinline{n}. Explain in words. Assume \lstinline{n} is a
non-negative integer and arithmetic and printing each take a fixed amount
of work.

\begin{center}
\begin{tabular}{ |p{13cm}|p{2cm}| }
\hline
\begin{lstlisting}
def mystery(n):
    total = 0
    for i in range(n // 2):
        j = 1
        while j < n:
            j *= 2
            total += i * j
    return total
\end{lstlisting} &  \\  \hline

\begin{lstlisting}
def foo(n):
    if n == 0:
        return
    for i in range(n // 2):
        print(i)
    foo(n // 2)
\end{lstlisting} &  \\  \hline
\end{tabular}
\end{center}

\end{blocksection}

\begin{solution}
\begin{itemize}
    \item \lstinline{mystery}: superlinear. The outer loop runs about half
    as many times as \lstinline{n}, and each iteration starts a fresh
    doubling loop. As \lstinline{n} grows, both the number of outer
    iterations and the work per iteration grow. Doubling \lstinline{n}
    roughly doubles the outer iterations and adds an inner iteration.
    \item \lstinline{foo}: linear. The loops do about half of \lstinline{n}
    steps, then a quarter, then an eighth, and so on. All those shrinking
    amounts together stay proportional to \lstinline{n}. The recursive
    calls do not each repeat the original amount of work.
\end{itemize}
\end{solution}
